\chapter{Stress Testing and Scenarios}\label{ch:rc:stress-testing-and-scenarios}

In March 2021 price falls in the most concentrated positions of Archegos Capital
Management, a family office, triggered margin calls it could not meet. It had grown in a year
from about 1.5 billion dollars of capital and 10 billion of exposure to more than 36 billion of
capital and 160 billion of exposure, much of it through total return swaps with several prime
brokers, each of which saw only its own part. Its default caused over 10 billion dollars of
losses across several large banks; one of them, Credit Suisse, lost close to 5.5 billion. The losses came from a
move that a risk measure calibrated on the recent past, applied to one broker's slice of the
book, would not have called likely. This chapter is about the tools that ask the other
question: not how much the book loses on a bad day, but what happens in a particular bad
world, which worlds would break the firm, and what supervisors require banks to survive.

\section{What value at risk cannot see}

\begin{definition}[Stress test]\label{def:rc:stress-testing-and-scenarios:stress}
A \emph{stress test}\index{stress test} revalues a portfolio, or a whole firm, under a
specified severe but plausible scenario of risk-factor moves, and reports the loss and its
consequences (capital, liquidity, limits), without attaching a probability to the scenario.
\end{definition}

VaR and ES summarise a distribution estimated from a window of history; they miss moves the
window does not contain, changes in correlation, the concentration and illiquidity of a
position (which lengthen the time to exit), and the losses of counterparties who hold the same
positions. A \omterm{def:rc:stress-testing-and-scenarios:stress}{stress test} replaces the distribution by a story, told in numbers.

\section{Historical scenarios}

\begin{definition}[Historical scenario]\label{def:rc:stress-testing-and-scenarios:historical}
A \emph{historical scenario}\index{historical scenario} applies to today's positions the moves
the \omterm{def:rc:market-risk-measures:factor}{risk factors} made over a past period: differences for rates and spreads, relative changes
for prices.
\end{definition}

\begin{example}[Three episodes on one book]\label{ex:rc:stress-testing-and-scenarios:historical}
Chapter~21's book (long ten-year Treasuries, short two-year, long euros, short yen, short a
EURUSD straddle), revalued in full under three periods of Treasury and ECB data:
\begin{itemize}
\item Lehman, 12 September to 10 October 2008: two-year yield $-61$ basis points, ten-year $+15$,
EURUSD $-3.52\%$, USDJPY $-7.85\%$; a loss of USD~17.26 million;
\item March 2020, 6 to 20 March: $-12$ and $+18$ basis points, EURUSD $-5.71\%$, USDJPY $+5.11\%$; a
loss of USD~20.13 million;
\item September 2022, 21 to 28 September: $+5$ and $+21$ basis points, EURUSD $-3.50\%$, USDJPY
$+0.46\%$; a loss of USD~11.61 million.
\end{itemize}
The ten-day 99\% historical VaR scaled from chapter~21 is USD~6.63 million: every episode loses two
to three times as much (\cref{fig:rc:stress-testing-and-scenarios:historical}). The short straddle,
harmless on normal days, dominates the losses in all three.
\end{example}

\begin{omfigure}
\begin{tikzpicture}
\begin{axis}[width=11.5cm,height=5.2cm,xbar,bar width=12pt,
  symbolic y coords={ten-day VaR,September 2022,March 2020,Lehman 2008},ytick=data,
  xlabel={loss (USD million)},
  tick label style={font=\small},label style={font=\small},
  xmin=0,xmax=24,grid=major,grid style={gray!25},
  nodes near coords,every node near coord/.append style={font=\footnotesize,/pgf/number format/fixed,/pgf/number format/precision=1},
  table/col sep=comma]
\addplot[fill=omThm!60,draw=omThm] table[y=scenario,x=loss]{figdata/rates-credit-risk/22-stress-testing-and-scenarios/historical.csv};
\end{axis}
\end{tikzpicture}

\omcaption{Losses of chapter~21's book under three \omterm{def:rc:stress-testing-and-scenarios:historical}{historical scenarios}, against its ten-day
99\% VaR. Each episode costs two to three times the VaR. Data: US Treasury par yields, ECB
reference rates; the chapter's tutorial.}\label{fig:rc:stress-testing-and-scenarios:historical}
\end{omfigure}

\begin{dated}{2026-09}{October 1987 and September 2022}\label{dat:rc:stress-testing-and-scenarios:history}
On 19 October 1987 the Dow Jones Industrial Average fell 508 points, 22.6\%, in one session, the
largest one-day decline in its history. In late September 2022 the thirty-year UK gilt yield rose
160 basis points in a few days, with an intraday range of 127 basis points on one day, as
liability-driven investment funds met margin calls.
\end{dated}

\section{Hypothetical scenarios and their consistency}

\begin{definition}[Hypothetical scenario]\label{def:rc:stress-testing-and-scenarios:hypothetical}
A \emph{hypothetical scenario}\index{hypothetical scenario} is a set of moves designed rather than
observed: a story (a rate spike, a devaluation, a default) translated into shocks to some \omterm{def:rc:market-risk-measures:factor}{risk
factors}, with the rest of the factors set consistently.
\end{definition}

\begin{definition}[Scenario conditioning]\label{def:rc:stress-testing-and-scenarios:conditioning}
\emph{Scenario conditioning}\index{scenario conditioning} fills the factors a scenario does not
specify with their expectation given the specified ones. Under a joint normal model with
covariance $\Sigma$, shocks $x_s$ to the specified factors give $\E[x_u\mid x_s] = \Sigma_{us}\Sigma_{ss}^{-1}x_s$ for the
others.
\end{definition}

\omcode{code/firm/stresstest/firm_stresstest.py}{32}{40}{Conditioning: the factors a scenario leaves open, at their conditional expectation.}\label{lst:rc:stress-testing-and-scenarios:condition}

\begin{example}[A ten-year shock]\label{ex:rc:stress-testing-and-scenarios:hypo}
Shock the ten-year yield up 100 basis points. Alone, it costs the book USD~14.80 million. Filled in
with the ten-day covariance of daily moves since 2020, the scenario also moves the two-year up 80
basis points, the euro down 0.36\% and USDJPY up 2.25\%; the loss is USD~10.04 million, because the
two-year short gains. The scenario sits 5.3 standard deviations from the centre of the ten-day
distribution.
\end{example}

\section{Reverse stress tests}

\begin{definition}[Reverse stress test]\label{def:rc:stress-testing-and-scenarios:reverse}
A \emph{reverse stress test}\index{reverse stress test} starts from an outcome, a loss that would
break a limit, the capital or the firm, and searches for the scenarios that produce it; its
quantitative form finds the most plausible one, the move of smallest Mahalanobis distance
$\sqrt{x^\top\Sigma^{-1}x}$ with loss at least $L$.
\end{definition}

\begin{proposition}[The most plausible loss scenario of a linear book]\label{prop:rc:stress-testing-and-scenarios:rev}
For P\&L $\delta^\top x$ and covariance $\Sigma$, the move of smallest Mahalanobis distance with loss $L$ is
$x^* = -L\,\Sigma\delta/(\delta^\top\Sigma\delta)$, at distance $L/\sqrt{\delta^\top\Sigma\delta}$: the loss measured in standard
deviations of the book's P\&L.
\end{proposition}
\begin{proof}
Minimise $\tfrac12x^\top\Sigma^{-1}x$ subject to $\delta^\top x = -L$: the Lagrange condition gives $\Sigma^{-1}x = \mu\delta$,
so $x = \mu\Sigma\delta$, and the constraint fixes $\mu$.
\end{proof}

For a non-linear book the method iterates: take the direction $\Sigma\nabla\mathrm{P\&L}$ at the current point,
move along it until the loss is reached, recompute the gradient there.

\begin{example}[Breaking a USD~10 million limit]\label{ex:rc:stress-testing-and-scenarios:rev}
With the ten-day covariance, the linear model's most plausible scenario for a USD~10 million loss
raises the two-year yield 33 basis points and the ten-year 58, lowers the euro 3.18\% and raises
USDJPY 1.78\%, at 3.93 standard deviations; revalued in full, that scenario loses USD~14.00
million. The full-revaluation search finds a nearer one, at 2.47 standard deviations: the euro
down 3.45\%, USDJPY up 2.15\%, yields up 10 and 17 basis points
(\cref{fig:rc:stress-testing-and-scenarios:reverse}). The short straddle makes large currency
moves in either direction the cheapest way to lose.
\end{example}

\begin{omfigure}
\begin{tikzpicture}
\begin{axis}[width=11.5cm,height=5.2cm,ybar,bar width=11pt,
  symbolic x coords={2y,10y,EURUSD,USDJPY},xtick=data,
  ylabel={move (ten-day standard deviations)},
  tick label style={font=\small},label style={font=\small},
  ymin=-3.5,ymax=3.5,grid=major,grid style={gray!25},
  legend columns=2,legend style={font=\footnotesize,at={(0.5,-0.2)},anchor=north,draw=none,fill=none,/tikz/every even column/.append style={column sep=8pt}},
  table/col sep=comma]
\addplot[fill=omDef!60,draw=omDef] table[x=factor,y=linear]{figdata/rates-credit-risk/22-stress-testing-and-scenarios/reverse.csv};
\addplot[fill=omThm!60,draw=omThm] table[x=factor,y=full]{figdata/rates-credit-risk/22-stress-testing-and-scenarios/reverse.csv};
\legend{linear model (3.93 sd),full revaluation (2.47 sd)}
\end{axis}
\end{tikzpicture}

\omcaption{The most plausible ten-day scenarios that cost the book USD~10 million, from the
linear model and from full revaluation, each factor's move in its own ten-day standard deviations.
The linear model spreads the move across rates and currencies; with the straddle's convexity, a
currency move does the damage at a much smaller distance. Data: the chapter's
tutorial.}\label{fig:rc:stress-testing-and-scenarios:reverse}
\end{omfigure}

\section{Supervisory stress tests}

\begin{definition}[Supervisory stress test]\label{def:rc:stress-testing-and-scenarios:supervisory}
A \emph{supervisory stress test}\index{supervisory stress test} is a \omterm{def:rc:stress-testing-and-scenarios:stress}{stress test} whose scenarios
are set by the supervisor and applied to many banks at once, to size capital buffers and compare
the banks' resilience.
\end{definition}

The US Federal Reserve's scenarios combine macroeconomic and financial paths with, for banks with
large trading operations, a global market shock applied to their trading and other fair-valued
positions and, for some, a counterparty-default component. The European Banking Authority runs
EU-wide tests on a common methodology.

\begin{dated}{2026-09}{Supervisory scenarios}\label{dat:rc:stress-testing-and-scenarios:fed}
The Federal Reserve's final 2026 severely adverse scenario (published in February 2026) is a severe
global recession triggered by a fall in risk appetite: the unemployment rate rises 5.5 percentage
points to a peak of 10\% in the third quarter of 2027, equity prices fall about 58\% through the
third quarter of 2026, house prices 30\% and commercial real estate prices 39\%, and the VIX peaks at
72. In the European Banking Authority's 2025 EU-wide \omterm{def:rc:stress-testing-and-scenarios:stress}{stress test} (64 banks, 75\% of EU banking assets) the
adverse scenario cost 370 basis points of common equity tier~1, leaving an aggregate ratio of 12\%; the
draft methodology of the 2027 exercise was published on 11 June 2026.
\end{dated}

\section{Tutorial: a scenario library}\label{tut:rc:stress-testing-and-scenarios:tutorial}

\begin{tutorial}
\textbf{Goal.} Apply \omterm{def:rc:stress-testing-and-scenarios:historical}{historical scenarios} to chapter~21's book, complete a \omterm{def:rc:stress-testing-and-scenarios:hypothetical}{hypothetical scenario} by
conditioning, and find the most plausible scenario behind a given loss. \textbf{End state:} the
numbers of
\cref{ex:rc:stress-testing-and-scenarios:historical,ex:rc:stress-testing-and-scenarios:hypo,ex:rc:stress-testing-and-scenarios:rev}
and the charts.
\begin{tutsteps}
\item \textbf{Library}: \texttt{scenarios()} builds the three episodes with \texttt{historical}.
\item \textbf{Run}: \texttt{historical\_table()}; compare with \texttt{var10()}.
\item \textbf{Condition}: \texttt{hypothetical(100)}.
\item \textbf{Reverse}: \texttt{reverse(10e6)}; \texttt{fig\_rc\_stress.py} writes the charts.
\end{tutsteps}
\textbf{What to change next.} Condition on the stressed covariance of March 2020 instead; set the
straddle to zero and repeat the \omterm{def:rc:stress-testing-and-scenarios:reverse}{reverse stress test}.
\end{tutorial}

\section{Build: the stress-testing engine}\label{bld:rc:stress-testing-and-scenarios:stresstest}

\begin{build}
\bfield{Purpose} The firm's scenario library and stress calculations, run nightly by the risk
engine of chapter~29 alongside VaR.
\bfield{Interface} \texttt{Scenario(name, moves)}; \texttt{historical(name, dates, levels,
start, end, log\_cols)}; \texttt{condition(cov, shocked, values)}; \texttt{mahalanobis};
\texttt{reverse\_linear(delta, cov, loss)}; \texttt{reverse\_nonlinear(pnl, cov, loss, x0)};
\texttt{run(pnl, scenarios)}.
\bfield{Rules} Moves in the model's units (yields in decimals, prices in log returns);
plausibility by Mahalanobis distance under the chosen covariance.
\bfield{Acceptance tests} \texttt{code/firm/stresstest/tests/}: historical moves; conditioning
equals the regression; the linear reverse scenario lies on the loss plane and is closer than any
other point on it; the non-linear search recovers it for a linear book.
\bfield{Stretch} Stressed covariances and fat-tailed plausibility; scenario libraries by
regime; liquidity-adjusted horizons per position; firm-wide reverse stress including funding.
\end{build}

\begin{omsources}
\item US Securities and Exchange Commission, press release 2022-70, 27 April 2022.
\item Paul, Weiss, report to the Special Committee of the Board of Credit Suisse Group on
Archegos Capital Management, 29 July 2021.
\item Federal Reserve History, ``Stock market crash of 1987''.
\item Board of Governors of the Federal Reserve System, \emph{2026 Stress Test Scenarios}, 2026.
\item Bank of England, letter from Sir Jon Cunliffe to the Treasury Committee, 5 October 2022.
\end{omsources}

\section{Exercises}

\begin{exercise}[$\star$]\label{exo:rc:stress-testing-and-scenarios:1}
Why can a stress loss be many times the VaR without the VaR model being wrong?
\end{exercise}

\begin{exercise}[$\star$]\label{exo:rc:stress-testing-and-scenarios:2}
In a two-factor normal model with volatilities 10 and 5 basis points and correlation 0.6, the
first factor is shocked by 40 basis points. What is the conditional move of the second?
\end{exercise}

\begin{exercise}[$\star$]\label{exo:rc:stress-testing-and-scenarios:3}
A linear book has a ten-day P\&L standard deviation of USD~2.5 million. At what Mahalanobis
distance does it lose USD~10 million?
\end{exercise}

\begin{exercise}[$\star\star$]\label{exo:rc:stress-testing-and-scenarios:4}
Why did the hypothetical ten-year shock cost less once completed by conditioning?
\end{exercise}

\begin{exercise}[$\star\star$]\label{exo:rc:stress-testing-and-scenarios:5}
Why is the linear reverse stress scenario not the most plausible one for this book?
\end{exercise}

\begin{exercise}[$\star\star$]\label{exo:rc:stress-testing-and-scenarios:6}
What does a \omterm{def:rc:stress-testing-and-scenarios:historical}{historical scenario} assume about today's positions and markets that may be false?
\end{exercise}

\begin{exercise}[$\star\star\star$]\label{exo:rc:stress-testing-and-scenarios:7}
\emph{Coding.} Run the \omterm{def:rc:stress-testing-and-scenarios:reverse}{reverse stress test} for a loss of USD~10 million without the straddle and
compare the distance.
\end{exercise}

\begin{exercise}[$\star\star\star$]\label{exo:rc:stress-testing-and-scenarios:8}
\emph{Find the flaw.} ``Our reverse stress scenario is 4 standard deviations away, so it happens
once in 30\,000 ten-day periods and we can ignore it.''
\end{exercise}

\section{Problem: The Family Office from the Broker's Side}

\begin{problem}[{Weekend problem --- when the margin runs out}]\label{pb:rc:stress-testing-and-scenarios:1}
A prime broker faces a client through total return swaps on five concentrated stocks, USD~20
billion of notional weighted 30, 25, 20, 15 and 10\%, hedged by holding the stocks. The client posts
margin of 7.5\% of notional. If the client defaults, the broker needs five days to sell. In an
illustrative one-factor model each stock has a beta of 1.2 to a market with 1.2\% daily volatility
and 3\% daily idiosyncratic volatility.

\medskip\noindent\textbf{Part I --- The book.}
\begin{enumerate}
\item Give the margin in dollars.
\item Give the standard deviation of the stocks' value over five days.
\item Why is the broker's risk the client's default, not the stocks?
\item What would a one-day 99\% VaR of the swaps say while the client pays its margin?
\item What did the broker not see that the client's other brokers held?
\end{enumerate}

\medskip\noindent\textbf{Part II --- Reverse stress.}
\begin{enumerate}[resume]
\item Give the most plausible five-day moves that consume the margin.
\item Give their Mahalanobis distance and the normal probability of a move at least that far.
\item Why do the largest positions fall the most in that scenario?
\item What uniform fall consumes the margin?
\item What changes with a margin of 20\%?
\end{enumerate}

\medskip\noindent\textbf{Part III --- Beyond the model.}
\begin{enumerate}[resume]
\item Why would five days be optimistic?
\item How do other brokers' sales change the scenario?
\item Why was the swap margin rate itself a risk?
\item What information would a dynamic margin require?
\item What single-name concentration limit would you set?
\end{enumerate}

\medskip\noindent\textbf{Part IV --- Judgement.}
\begin{enumerate}[resume]
\item Why is this a stress-testing problem and not a VaR problem?
\item How should the reverse stress result enter the broker's decisions?
\item What would you ask the client before extending more capacity?
\item State the \emph{named result}: the moves at which the book exhausts its margin, and their
distance.
\item In one sentence: what does a \omterm{def:rc:stress-testing-and-scenarios:reverse}{reverse stress test} tell a risk committee?
\end{enumerate}
\end{problem}

\section{Interview questions}

\begin{interviewq}[$\star$ \iqroles{risk}]\label{iq:rc:stress-testing-and-scenarios:1}
What is the difference between VaR and a \omterm{def:rc:stress-testing-and-scenarios:stress}{stress test}? When would you trust each?
\end{interviewq}

\begin{interviewq}[$\star\star$ \iqroles{risk, researcher}]\label{iq:rc:stress-testing-and-scenarios:2}
How do you make a \omterm{def:rc:stress-testing-and-scenarios:hypothetical}{hypothetical scenario} internally consistent?
\end{interviewq}

\begin{interviewq}[$\star\star$ \iqroles{researcher}]\label{iq:rc:stress-testing-and-scenarios:3}
Derive the most plausible scenario for a given loss of a linear portfolio.
\end{interviewq}

\begin{interviewq}[$\star\star$ \iqroles{bank, risk}]\label{iq:rc:stress-testing-and-scenarios:4}
What does the Federal Reserve's global market shock test, and how does it differ from VaR?
\end{interviewq}

\begin{interviewq}[$\star\star\star$ \iqroles{risk, trader}]\label{iq:rc:stress-testing-and-scenarios:5}
How would you \omterm{def:rc:stress-testing-and-scenarios:stress}{stress test} a prime brokerage book?
\end{interviewq}

\begin{interviewq}[$\star\star\star$ \iqroles{developer}]\label{iq:rc:stress-testing-and-scenarios:6}
Design a scenario library for a bank's nightly stress run. What goes in it, and how is it
governed?
\end{interviewq}
